\section{Ejercicio VII}

El CVA (\textit{credit valuation adjustment}) es un ajuste contable que reduce el valor de un derivado financiero para reflejar el riesgo de que la contraparte incumpla sus obligaciones. Técnicamente, representa el valor presente del costo esperado que asumiría una institución si la otra parte del contrato quiebra antes de liquidar la operación.

La necesidad de calcular y reportar el CVA nace tras la crisis financiera de 2008. Durante ese periodo, las instituciones registraron pérdidas multimillonarias, no por quiebras directas, sino por el deterioro en la calidad crediticia de sus contrapartes, lo que disminuyó el valor de mercado de sus portafolios. Esto impulsó al Comité de Basilea (bajo el marco de Basilea III) a exigir que las instituciones midieran este riesgo y mantuvieran capital suficiente para respaldarlo.

Se aplica en la valoración de instrumentos derivados negociados bilateralmente en mercados extrabursátiles. Dado que estas operaciones no pasan por una cámara de compensación central que actúe como garante, ambas partes asumen un riesgo de crédito mutuo que debe ser cuantificado.

El CVA impacta los estados financieros en dos frentes:

Balance General: El activo financiero se presenta por su valor neto, es decir, el valor del derivado libre de riesgo menos el monto del ajuste por CVA.

Estado de Resultados: Las fluctuaciones del riesgo crediticio son dinámicas. Si la probabilidad de incumplimiento de la contraparte aumenta, el CVA crece y este incremento se reconoce inmediatamente como una pérdida en el periodo contable, reduciendo la utilidad neta.